\chapter{Tier 2B: forecasting and chronological operation}\label{chap:tier2}
\section{Information and experimental design}\label{sec:design2}
The operating decision is made once per day. Following the classroom
simplification, current-day generation, demand and price are known at that
decision, while future days are unknown \cite{teacher}. Let
\(y_t=(G_t,D_t,p_t)\) and let \(\mathcal I_t\) contain observations through
day \(t\), the installed equipment and the previous day's usable stock.
Only information in \(\mathcal I_t\) can determine the implemented action.
Current-day aggregates are treated as observed at this daily planning decision.

We reserve 2020--22 for training and capacity selection, 2023 for model and
horizon selection, and 2024--25 for the final test. A forecasting training
target must lie before its corresponding cutoff. Repair calibration, scaling and hyperparameter
selection use training/validation data only. The six-year Tier 1 capacity
comparison remains a hindsight benchmark; test observations must not
choose the capacity used in the predictive experiment.

\section{Multi-step forecasting model}\label{sec:forecast2}
The proposed first predictive model is regularised linear autoregression,
compared against a seven-day seasonal-naive predictor. At each forecast
origin \(t\), the feature vector \(f_t\) contains observed generation,
demand and price lags, day-of-week indicators and annual sine/cosine terms.
For variable \(j\in\{G,D,p\}\) and horizon \(k=1,\ldots,H\), fit
\begin{equation}
\widehat\beta_{jk}=\arg\min_{\beta}
\sum_{\tau\in\mathcal A_k}(y_{\tau+k,j}-f_\tau^{\top}\beta)^2
+\lambda\lVert\beta\rVert_2^2,
\qquad
\widehat y_{t+k,j}=f_t^{\top}\widehat\beta_{jk},
\end{equation}
where \(\mathcal A_k\) contains eligible training origins whose targets
are already observed, \(\lambda\ge0\) is the validation-selected penalty
and \(H\) is the planning horizon. Forecast generation and demand are
projected onto nonnegative values; negative price forecasts remain allowed.
The seasonal-naive forecast uses the most recent observation with the same
weekday, recursively if needed.

\section{Rolling optimisation and storage feedback}\label{sec:rolling2}
At day \(t\), installed battery and electrolyser counts are fixed. A
planning problem uses observed values for the current day and predictions
for \(t+1,\ldots,t+H\). It retains the Tier 1 flow, demand, throughput and electrolysis
constraints, but replaces the pooled energy balance with chronological stock:
\begin{equation}
s_\tau=s_{\tau-1}+\eta c_\tau-r_\tau/\eta,
\qquad 0\le s_\tau\le E_Bb,
\qquad \tau=t,\ldots,t+H.
\end{equation}
Here \(s_{t-1}\) is the measured initial stock for that solve. The optimiser
maximises forecast operating receipts plus a terminal-stock value; the latter
is selected on validation data to avoid repeatedly liquidating storage at an
artificial short horizon. Only \((q_t,c_t,r_t,h_t)\) is implemented. Its
realised stock becomes the next solve's initial condition. The test's initial
and final stock treatment is held equal across all compared policies.

\section{Evaluation and results to be completed}\label{sec:evaluation2}
Forecast error is reported by series and horizon using mean absolute error
and root mean square error, with price errors in \gbp/MWh and physical
errors in MWh. Operating evaluation compares the predictive controller,
the same controller with seasonal-naive forecasts, a current-day rule and
a full-path perfect-information optimiser, all at identical installed
capacities and with identical data, costs and terminal stock. The gap to
perfect information measures the realised loss from restricted information,
rather than a change in investment scale.

Since zero batteries remove cross-day storage, fixed nonzero battery cases
will diagnose prediction and feedback value alongside the training-selected
investment. Results will include forecast errors, realised annual value,
surplus capture, lost sales, stock feasibility and seasonal/adverse-period
comparisons. These forecast and control experiments are pending.
